\documentclass[letterpaper]{article} % DO NOT CHANGE THIS
\usepackage[submission]{aaai2027}  % DO NOT CHANGE THIS
\usepackage[hyphens]{url}  % DO NOT CHANGE THIS
\usepackage{graphicx} % DO NOT CHANGE THIS
\urlstyle{rm} % DO NOT CHANGE THIS
\def\UrlFont{\rm}  % DO NOT CHANGE THIS
\usepackage{natbib}  % DO NOT CHANGE THIS AND DO NOT ADD ANY OPTIONS TO IT
\usepackage{caption} % DO NOT CHANGE THIS AND DO NOT ADD ANY OPTIONS TO IT
\frenchspacing  % DO NOT CHANGE THIS
\usepackage[utf8]{inputenc} % allow utf-8 input
\usepackage[T1]{fontenc}    % use 8-bit T1 fonts
\usepackage{booktabs}       % professional-quality tables
\usepackage{multirow}       % multi-row cells for compact tables
\usepackage{amsfonts}       % blackboard math symbols
\usepackage{nicefrac}       % compact symbols for 1/2, etc.
\usepackage{microtype}      % microtypography
\usepackage{xcolor}         % colors
\usepackage{amsmath, amsthm, amssymb}
\usepackage{mathtools}
\usepackage{algorithm, algorithmic}
\usepackage{tikz, pgfplots}
\usetikzlibrary{shapes.geometric, arrows.meta, positioning, fit, backgrounds}

\newtheorem{theorem}{Theorem}
\newtheorem{definition}{Definition}
\newtheorem{corollary}{Corollary}
\newtheorem{proposition}{Proposition}
\newtheorem{remark}{Remark}
\newtheorem{assumption}{Assumption}

\newcommand{\E}{\mathbb{E}}
\newcommand{\GOMDP}{\text{GOMDP}}
\newcommand{\Conf}{\text{Conf}}
\newcommand{\HA}{\text{HA}}
\newcommand{\Pbreak}{P_{\text{break}}}
\newcommand{\Le}{L_e}

\usepgfplotslibrary{groupplots}
\pgfplotsset{compat=1.18}
\pdfinfo{
/TemplateVersion (2027.1)
}

\setcounter{secnumdepth}{2}

% --- Macros ---
\newcommand{\Rset}{\mathbb{R}}
\newcommand{\calA}{\mathcal{A}}
\newcommand{\calS}{\mathcal{S}}
\newcommand{\calT}{\mathcal{T}}
\newcommand{\calG}{\mathcal{G}}
\newcommand{\calU}{\mathcal{U}}
\newcommand{\calB}{\mathcal{B}}
\newcommand{\calE}{\mathcal{E}}
\newcommand{\calX}{\mathcal{X}}
\newcommand{\calZ}{\mathcal{Z}}

% --- Appendix pointers ---
% The appendices are the separately compiled supplementary document, so their
% labels are not visible to this file. These letters mirror the section order
% in supplementary.tex; update both together.
\newcommand{\appNotation}{A}
\newcommand{\appProof}{B}
\newcommand{\appProofTwo}{C}
\newcommand{\appByzantine}{D}
\newcommand{\appSensitivity}{E}
\newcommand{\appMultisig}{F}
\newcommand{\appVIIRS}{G}

\title{Governance-Invariant MDPs: A Framework and Formal Safety Case for Agentic Wildfire Monitoring}
\author{Anonymous Submission}
\affiliations{Anonymous Affiliation}

\begin{document}

\maketitle

\begin{abstract}
Wildfire early warning must be fast enough to save lives yet accountable enough to trust; the authority to broadcast a public alert cannot be ceded to an unaccountable autonomous agent. Constrained reinforcement learning (RL) enforces safety only in expectation and therefore permits violations. We introduce the \emph{Governance-Invariant MDP} (GOMDP), which relocates mandatory human authorization from a policy penalty to a cryptographic state-transition invariant in the environment. Theorem~1 provides a policy-agnostic safety guarantee; Corollary~1 states the composite bound governing deployment, whose residual risk is dominated not by signature strength but by validator compromise, bounded in closed form by Theorem~2 and extended to correlated compromise by Proposition~1. Instantiated as a Python reference implementation of Hyperledger Fabric chaincode logic and evaluated on three VIIRS-derived fire events, PPO-GOMDP reduces the false public alert rate from 22.4\% to 6.0\% at full authorization compliance, with latency statistically equivalent to constrained RL (TOST, $p=0.004$). A layered ablation prices each enforcement layer: Ed25519 verification alone blocks 100/100 injection attacks, Byzantine consensus adds nothing measurable, and under correlated compromise its analytic advantage over a hardened single verifier disappears. We also report the cost that the accountability literature routinely omits: the gate raises missed detections from 0.9\% to 2.1\%, and to 7.8\% under operator fatigue.
\end{abstract}

\section{Introduction}
\label{sec:intro}

Wildfire early warning, under growing pressure from climate-driven extremes~\citep{divirgilio2019}, is shifting from human monitoring toward autonomous, agentic AI that fuses satellite, UAV, and sensor data to decide when to raise a public alarm. The stakes are asymmetric: a false alert can trigger a costly and dangerous mass evacuation, while a missed detection can cost lives, so the decision to broadcast must remain accountable to a named human authority and auditable after the fact. Constrained MDPs~\citep{altman1999} and trust-region methods such as CPO~\citep{achiam2017} enforce constraints only in expectation via Lagrangian relaxation, so violations persist after convergence. Projection-based updates~\citep{yang2020pcpo}, distributional critics~\citep{yang2021wcsac}, and safety layers~\citep{dalal2018} reduce but cannot eliminate them, because all tie the guarantee to policy or model quality~\citep{gu2022survey}. For public warning, an in-expectation, unaccountable guarantee is the wrong instrument.

We address this gap with the \emph{Governance-Invariant MDP} (GOMDP), in which the safety constraint resides in the environment's transition kernel under cryptographic enforcement rather than in the policy's reward signal. The structural insight is a separation of concerns familiar from enforcement theory~\citep{schneider2000,ligatti2005}: safety of this kind is an environment property, not a policy property. Any policy---random, greedy, or fully trained---satisfies the governance constraint, because the constraint is a state-transition invariant that no action can bypass. The guarantee is deliberately narrow: it secures \emph{authorization integrity and accountability}, meaning no alert without a valid and non-repudiable human authorization, and says nothing about whether the human decided correctly, which no cryptography can supply.

Two commitments shape how we report results. We state the guarantee at the strength it actually holds: Theorem~\ref{thm:safety} is conditional on faithful implementation and at most $f$ Byzantine validators, and Corollary~\ref{cor:composite} makes the composite bound explicit, dominated by validator compromise rather than signature forgery. We also report what the mechanism costs as carefully as what it buys, including a negative result about our own architecture.

\noindent\textbf{Contributions.} \textbf{(C1)}~\emph{A well-typed formalism for environment-level enforcement}: we define the GOMDP over a governance-observable state space so that the governance predicate is a genuine function of $(s,a)$ and the constrained kernel is stationary and total, and show that the resulting safety certificate is independent of the policy's optimality gap (Corollary~\ref{cor:decoupling}), unlike CMDP, CPO, and WCSAC. \textbf{(C2)}~\emph{An honestly scoped safety case}: Theorem~\ref{thm:safety} identifies the minimal assumptions; Corollary~\ref{cor:composite} states the composite bound those assumptions leave; Theorem~\ref{thm:adversarial} bounds validator compromise in closed form; and Proposition~\ref{prop:correlated} extends the bound to correlated compromise, where the case for consensus over a hardened single verifier collapses. \textbf{(C3)}~\emph{A layered decomposition with a negative result}: separating authentication from consensus (the Central+Sig ablation) shows that signature verification alone defeats injection and that consensus adds nothing measurable---a decomposition absent from prior shielding and safety-layer evaluations. \textbf{(C4)}~\emph{A full cost ledger}: governance adds $4.2$ steps to the detection-to-broadcast path and reduces false public alerts from 22.4\% to 6.0\%, but raises missed detections $2.3\times$ and up to $8.7\times$ under operator fatigue.

\section{Related Work}
\label{sec:related}

\paragraph{Constrained RL and safe AI.} Altman's CMDP framework~\citep{altman1999} enforces constraints via Lagrangian relaxation with in-expectation guarantees; CPO~\citep{achiam2017} and RCPO~\citep{tessler2019} reduce violations without eliminating them, and PCPO~\citep{yang2020pcpo} projects updates onto the constraint set. WCSAC~\citep{yang2021wcsac} adds a distributional safety critic but remains policy-level, and \citet{dalal2018} project actions through a learned linearized cost model that inherits model error. Safe shielding~\citep{alshiekh2018} filters unsafe actions at the environment boundary and is architecturally closest to GOMDP; our shield baseline shows that this enforcement is complete against the policy's own actions yet defenseless against external injection. Surveys~\citep{gu2022survey} confirm that per-trajectory guarantees under adversarial conditions remain open for these methods.

\paragraph{Enforcement theory, runtime assurance, and TEEs.} \citet{schneider2000} characterize which security policies an execution monitor can enforce, and \citet{ligatti2005} extend this with edit automata. The GOMDP predicate is a safety property in exactly this sense and $\calT_G$ is a truncation automaton over alert actions; what we add is a monitor that is cryptographically authenticated rather than trusted by assumption. The Simplex architecture~\citep{sha2001} pioneered environment-boundary enforcement in aerospace, and GOMDP inherits its philosophy while replacing the trusted local monitor with a distributed one, shifting the threat model from component failure to adversarial compromise. Control barrier functions~\citep{ames2019} enforce forward-invariant safe sets for continuous dynamics, whereas our predicate is discrete and requires non-repudiation, and runtime monitoring~\citep{leucker2009} observes events rather than embedding enforcement in the kernel. A trusted execution environment (TEE)~\citep{costan2016} is the obvious lower-complexity alternative, supplying tamper-evidence and remote verifiability without a quorum; what it does not supply is attributable multi-party authorization across mutually distrusting agencies, which is the residual case for consensus.

\paragraph{Wildfire detection and blockchain accountability.} Contemporary systems apply deep learning and multimodal fusion~\citep{ghali2022}, and IoT threshold pipelines~\citep{tarigan2026} report detection latency near 45 time units without adaptive patrol. \citet{partheepan2023} identify governance and accountability as open challenges in autonomous UAV monitoring, and \citet{gao2024} address a related domain with soft Lagrangian constraints. Permissioned blockchains such as Hyperledger Fabric~\citep{androulaki2018} enable authenticated, tamper-evident execution under BFT ordering, while human-in-the-loop (HITL) architectures generally treat oversight as an advisory or interface-level check~\citep{holzinger2025}. Where prior disaster-monitoring work has used ledgers, it has more often been for post-hoc auditing than for gating the alert transition itself, \citet{akarma2026} being the exception we know of; Assumption~\ref{ass:impl} connects our guarantee to smart-contract verification~\citep{bhargavan2016}.

\section{The Governance-Invariant MDP}
\label{sec:gomdp}

\subsection{Definition}
\label{sec:def}

The construction requires the governance predicate to be a genuine function of the state and action, which in turn requires the state to carry the verification statistic and the authorization register explicitly.

\begin{definition}[Governance-Invariant MDP]
\label{def:gomdp}
A $\GOMDP$ is a tuple $(\calS, \calA, \calA_{\mathrm{alert}}, a_\emptyset, \calT, \calT_G, r, \Omega, O, \calG)$ in which $\calA$, $r$, $\Omega$, $O$ are standard POMDP components with $r:\calS\times\calA\to\Rset$, and $\calT:\calS\times\calA\to\Delta(\calS)$ is the unconstrained transition kernel. The state space is \emph{governance-observable}: $\calS=\calX\times\calZ$, where $x\in\calX$ is the environment state and $z=(c,h)\in\calZ$ records the verification statistic $c\in[0,1]$ and the authorization register $h\in\{0,1\}$, both written by $\calT$ because the verification module and the human authorizer are components of the environment. The governance predicate $\calG:\calS\times\calA\to\{0,1\}$ depends on $s=(x,z)$ only through $z$. Writing $\calA_{\mathrm{alert}}\subset\calA$ for the designated alert actions and $a_\emptyset\in\calA\setminus\calA_{\mathrm{alert}}$ for a distinguished no-op, the governance-constrained kernel is
\begin{equation}
\calT_G(\cdot \mid s,a) =
\begin{cases}
\calT(\cdot \mid s,a), & a\notin\calA_{\mathrm{alert}} \text{ or } \calG(s,a)=1,\\[2pt]
\calT(\cdot \mid s,a_\emptyset), & a\in\calA_{\mathrm{alert}},\ \calG(s,a)=0.
\end{cases}
\label{eq:kernel}
\end{equation}
Enforcement of $\calG$ is \emph{cryptographic}: the environment executes an alert transition only when presented with a valid governance certificate, independently of the policy.
\end{definition}

Blocking redirects to the no-op branch rather than transitioning outside $\calS$, so $\calT_G$ is stationary and total on $\calS\times\calA$ by construction, and no alert-flagged successor is reachable when $\calG=0$. The distinction from a CMDP is structural: a CMDP constraint modifies the learning signal through a Lagrange multiplier, so a sufficiently suboptimal or adversarial policy can still act unsafely, whereas a $\GOMDP$ environment cannot produce an alert-flagged successor unless $\calG(s,a)=1$. The same separation distinguishes GOMDP from logical shielding~\citep{alshiekh2018}, which lacks non-repudiation, and from the safety layer of \citet{dalal2018}, whose learned projection inherits model error.

\begin{assumption}[Faithful Implementation]
\label{ass:impl}
The deployed chaincode correctly realizes $\calT_G$ as specified in Definition~\ref{def:gomdp}: for every $(s,a)$ with $a\in\calA_{\mathrm{alert}}$ it evaluates $\calG(s,a)$ and redirects to the no-op branch whenever $\calG(s,a)=0$. Formal verification of the chaincode, for which established techniques exist~\citep{bhargavan2016}, is outside our scope; Section~\ref{sec:discussion} returns to the consequences.
\end{assumption}

\subsection{Policy-Agnostic Safety}
\label{sec:theorem1}

\begin{theorem}[Policy-Agnostic Safety]
\label{thm:safety}
Let $\Sigma=(\GOMDP,\Pi_{\mathrm{Ed25519}},\mathcal{V})$ satisfy Assumption~\ref{ass:impl}, where $\Pi_{\mathrm{Ed25519}}$ is a $(t_{\mathrm{adv}},\epsilon_{\mathrm{sig}})$-secure signature scheme~\citep{bernstein2012ed25519} and $\mathcal{V}$ is a set of $k$ validators of which at most $f$ are Byzantine, with $k\geq 3f+1$. Let $\pi$ be any stochastic policy, including arbitrarily suboptimal or adversarial ones, operating within $\Sigma$ under $\calG_t\coloneqq\calG(s_t,a_t)=\mathbf{1}[\,c_t>\tau \wedge h_t=1\,]$. Then for any adversary running in time $t\leq t_{\mathrm{adv}}$ and all horizons $T$,
\begin{equation}
P_\pi\!\left(\exists\, t \leq T:\ \mathrm{Alert}_t = 1 \wedge \calG(s_t,a_t)=0\right) \leq T\,\epsilon_{\mathrm{sig}},
\label{eq:safety}
\end{equation}
which is negligible for Ed25519 ($\epsilon_{\mathrm{sig}}\approx 2^{-128}$). Conditional on Assumption~\ref{ass:impl} and on at most $f$ Byzantine validators, the guarantee is independent of $\pi$.
\end{theorem}

\begin{proof}[Proof sketch]
By reduction. Under Assumption~\ref{ass:impl} an adversary producing an unauthorized alert must either forge a valid certificate without an authorized validator key, which yields an Ed25519 forger with the same success probability and is bounded by $\epsilon_{\mathrm{sig}}$, or corrupt more than $f$ validators, which the hypothesis excludes. A union bound over $T$ timesteps gives~\eqref{eq:safety}. The full proof is in Appendix~\appProof{} of the supplementary material, to which all appendix references below refer.
\end{proof}

Theorem~\ref{thm:safety} is conditional in two places, and both conditions can fail in deployment: Assumption~\ref{ass:impl} is discharged only by formal verification, which we have not performed, and the Byzantine hypothesis is a probabilistic event whose failure probability Theorem~\ref{thm:adversarial} quantifies. Omitting either would misdescribe the guarantee, so we state the composite bound and use it, rather than~\eqref{eq:safety}, wherever the paper makes a deployment claim.

\begin{corollary}[Composite Safety Bound]
\label{cor:composite}
Under Assumption~\ref{ass:impl}, with validator compromise treated as a random event rather than assumed bounded by $f$, and with independent per-validator compromise probability $p_c$,
\begin{equation}
P_\pi\!\left(\exists\, t\leq T:\ \mathrm{Alert}_t=1 \wedge \calG_t=0\right) \leq T\,\epsilon_{\mathrm{sig}} + \Pbreak^{\GOMDP},
\label{eq:composite}
\end{equation}
with $\Pbreak^{\GOMDP}$ given by Eq.~\eqref{eq:pbreak_gomdp}. For $T=3000$ and $\epsilon_{\mathrm{sig}}=2^{-128}$ the first term is $\approx 9\times 10^{-36}$, so the bound is dominated by the second: $3.8\times 10^{-3}$ at $p_c=0.05$ and $0.35$ at $p_c=0.30$ for $k=7$, $f=2$. Residual risk is therefore governed by validator hygiene, not cryptographic strength.
\end{corollary}

\begin{corollary}[Safety--Optimality Decoupling]
\label{cor:decoupling}
Bounds~\eqref{eq:safety} and~\eqref{eq:composite} hold for any policy $\pi$, including the worst-case policy that minimizes return, so the optimality gap does not enter the safety certificate. This contrasts with CMDP, CPO, and WCSAC, where constraint violation probability depends directly on the optimality gap during training~\citep{achiam2017,yang2021wcsac}.
\end{corollary}

\begin{remark}[Novelty and scope]
\label{rem:scope}
That an environment blocking unauthorized transitions by construction cannot be made to permit them follows from Definition~\ref{def:gomdp}, and that secure signatures resist forgery is standard. The contribution of Theorem~\ref{thm:safety} is the connective step: it identifies the minimal assumptions under which environment-level enforcement transfers a guarantee to \emph{any} policy, and its reduction locates exactly where the guarantee fails---a broken primitive, a validator supermajority, or an implementation bug. Corollary~\ref{cor:composite} prices the second, which dominates.
\end{remark}

\subsection{Adversarial Robustness}
\label{sec:theorem2}

\begin{theorem}[Adversarial Robustness]
\label{thm:adversarial}
Under the threat model of Section~\ref{sec:threat}, with $k$ validators, tolerance $f\leq\lfloor (k-1)/3\rfloor$, and independent per-validator compromise probability $p_c$, the probability of a successful alert-injection breach in a $\GOMDP$ is bounded by the probability that more than $f$ validators are compromised:
\begin{equation}
\Pbreak^{\GOMDP} \leq \sum_{i=f+1}^{k}\binom{k}{i}p_c^{\,i}(1-p_c)^{k-i}.
\label{eq:pbreak_gomdp}
\end{equation}
For a centralized signature-verifying server, breach requires compromising the single verifier, so $\Pbreak^{\mathrm{sig}}=p_c$. Without authentication, breach requires only channel access, so $\Pbreak^{\mathrm{central}}=p_{\mathrm{attack}}$. A full proof is in Appendix~\appProofTwo.
\end{theorem}

For $k=7$, $f=2$, the exact tails are $0.004$, $0.026$, $0.148$, and $0.353$ at $p_c=0.05, 0.10, 0.20, 0.30$; the last is a pessimistic stress point, not a defensible operating regime. The $k=7$ tail falls below $p_c$ only for $p_c\lesssim 0.256$, above which a single verifier is preferable. Independence in~\eqref{eq:pbreak_gomdp} is also the assumption under which consensus looks best, so we relax it.

\begin{proposition}[Correlated Compromise]
\label{prop:correlated}
Let validator compromise events be exchangeable with marginal probability $p_c$ and intra-class correlation $\rho\in[0,1)$, modeled as $X\sim\mathrm{BetaBin}(k,\alpha,\beta)$ with $\alpha=p_c(1-\rho)/\rho$ and $\beta=(1-p_c)(1-\rho)/\rho$, so that $\Pbreak^{\GOMDP}=P(X>f)$. For $k=7$, $f=2$, moving from $\rho=0$ to $\rho=0.3$ raises the bound from $0.004$ to $0.048$ at $p_c=0.05$ and from $0.026$ to $0.103$ at $p_c=0.10$, and it moves the crossover against an equally exposed single verifier from $p_c\approx 0.256$ down to $p_c\approx 0.078$. If the single verifier is genuinely hardened, so that its exposure is $p_c^{\mathrm{sig}}=\tfrac{1}{2}p_c$, the crossover falls to $p_c\approx 0.152$ under independence and \emph{no crossover exists} for $\rho\geq 0.2$: consensus never dominates.
\end{proposition}

Proposition~\ref{prop:correlated} states when the consensus layer is analytically worth its cost. Shared software stacks, key-management practices, and co-located operators characterize realistic deployments and are exactly the conditions under which Byzantine agreement degenerates. Section~\ref{sec:ablation} reaches the same conclusion empirically.

\section{Instantiation for Wildfire Monitoring}
\label{sec:instantiation}

\subsection{Architecture and Threat Model}
\label{sec:arch}

Figure~\ref{fig:architecture} shows the pipeline. A heterogeneous sensor ensemble (UAV thermal cameras with detection probability $0.85$, ground IoT sensors, geostationary infrared feeds, and meteorological stations) generates observations fused probabilistically~\citep{khaleghi2013}; the VIIRS evaluation replaces this layer with preprocessed 375\,m active-fire data~\citep{schroeder2014}. A wildfire-risk digital twin drives fire-risk prediction, UAV coordination, and two-stage anomaly verification, with EMA-based threshold adaptation between episodes. Anomaly events are committed under BFT ordering ($k=7$, $f=2$), and the contract atomically checks the confidence threshold and the Ed25519 validator signature before dissemination. \textbf{Nothing is deployed}: the signing, key authorization, nonce ledger, and threshold logic are a Python reference implementation of the chaincode logic; BFT ordering is a delay distribution; and no UAV or Fabric network was instantiated.

\begin{figure*}[!t]
\centering
\includegraphics[width=0.80\textwidth]{images/Arch.pdf}
\caption{The wildfire GOMDP instantiation. The contract enforces $\calG$~\eqref{eq:governance} as a state-transition invariant, so every policy satisfies Theorem~\ref{thm:safety}. Red panels mark the threat model. No physical UAVs and no blockchain network were deployed.}
\label{fig:architecture}
\end{figure*}

\paragraph{Threat model.}
\label{sec:threat}
We consider five adversarial capabilities, all evaluated in Section~\ref{sec:robustness}: \textbf{sensor spoofing} (fabricated heat readings with probability $p_{\mathrm{spoof}}$, i.i.d.\ or targeted at high-belief cells~\citep{shafique2021}); \textbf{data tampering}~\citep{huang2021}; \textbf{alert injection}~\citep{he2015}; \textbf{validator compromise}; and \textbf{denial of service} (packet drops $p_{\mathrm{drop}}$~\citep{duo2022}). GOMDP secures authorization integrity, not decision correctness: a validator may authorize incorrectly, and on-chain logging makes such decisions attributable without preventing them.

\subsection{State, Objective, and Metrics}
\label{sec:objective}

The wildfire $\GOMDP$ has environment state $x_t=(H_t,W_t,D_t,R_t)$ (heat, weather, UAV configuration, risk estimate) and governance component $z_t=(c_t,h_t)$, with belief $b_t=P(x_t\mid o_{1:t},a_{1:t-1})$ updated by Bayesian filtering. The governance predicate is
\begin{equation}
\calG(s_t,a_t)=\mathbf{1}\!\left[\,c_t>\tau \ \wedge\ h_t=1\,\right], \qquad \tau=\tau_2=0.80,
\label{eq:governance}
\end{equation}
where $c_t=\Conf_t^{(2)}$ is the stage-2 confidence written into the state by the verification module and $h_t=\HA_t$ is the authorization register written by the operator.

\paragraph{Metrics.} Let $\calE$ be the episode's ignition events and $\calB$ its alert broadcasts, with $A(b)$ the geofenced footprint of $b$ at time $t(b)$. Detection latency is $L_d=t_{\mathrm{det}}-t_{\mathrm{ign}}$, where $t_{\mathrm{det}}$ is the first step at which $c_t>\tau_2$ for a cell inside the burning region, censored at $T$. The false public alert and missed-detection rates are false-discovery and miss ratios,
\begin{equation}
\begin{aligned}
F_p&=\frac{|\{b\in\calB: A(b) \text{ has no burning cell at } t(b)\}|}{|\calB|},\\
FN_r&=\frac{|\{e\in\calE: \nexists\, b\in\calB,\, e\in A(b)\}|}{|\calE|}.
\end{aligned}
\label{eq:fp}
\end{equation}
End-to-end latency is $\Le=L_d+\Delta_{\mathrm{BC}}+\Delta_{\mathrm{HITL}}$. The per-step quantities in~\eqref{eq:objective} are the indicators accumulating to these episode metrics: $L_{d,t}=\mathbf{1}[\text{ignition active and undetected at } t]$, so $\sum_t L_{d,t}=L_d$, and $F_{p,t}=\mathbf{1}[\text{broadcast at } t \text{ unwarranted}]$, so $\sum_t F_{p,t}=F_p|\calB|$. Injections mounted in Section~\ref{sec:robustness} do not enter $\calB$, so benign-episode $F_p$ and injection counts are distinct measurements.

The optimization objective is a scalar cost to be minimized,
\begin{equation}
J(\pi)=\E_\pi\!\left[\textstyle\sum_{t=1}^{T} c^{\mathrm{cost}}_t\right], \quad c^{\mathrm{cost}}_t=\alpha\tilde{L}_{d,t}+\beta\tilde{F}_{p,t}+\gamma\tilde{C}_{r,t},
\label{eq:objective}
\end{equation}
with $\alpha+\beta+\gamma=1$ and $C_{r,t}=\sum_{u\in\calU} d(u,p^*_u)/N$ the normalized repositioning cost. Each contribution is min-max normalized to $[0,1]$ using running extrema over the training trajectory, computed on validation seeds 100--109 disjoint from evaluation seeds 0--19, before weights $(0.50,0.35,0.15)$ apply; extrema stabilize within 50 episodes.

\subsection{Policies, Verification, and Contract}
\label{sec:policies}

\paragraph{PPO-GOMDP.} We train a PPO agent~\citep{schulman2017} inside the governed environment. The constraint resides in the kernel rather than the reward, and Theorem~\ref{thm:safety} holds before the agent learns anything. The policy is a two-layer MLP (hidden dimensions 256 and 128, ReLU) mapping a normalized flattened belief vector $\hat{b}_t\in\Rset^{10000+2N}$ to sector-assignment logits. The joint action space over $N=20$ UAVs and $Z=25$ sectors is factorized as $N$ conditionally independent categoricals sharing the trunk, one head per UAV, with sampled duplicates resolved by the proximity rule of Algorithm~\ref{alg:coord}; this avoids autoregressive decoding over $25^{20}$ assignments at the cost of ignoring within-step coupling. Alert generation belongs to the verification pipeline, not the policy: the policy controls where to look; the predicate controls whether to alert. Training runs 1000 episodes at learning rate $3\times10^{-4}$, clip ratio $0.2$, entropy bonus $0.01$, discount $0.99$, 4 update epochs, and batch size 64. Returns are negative, so we rescale by $\max_{e\leq 50}|G_e|$ over the first 50 episodes into $[-1,0]$. Training stops when validation $L_d$ fails to improve by more than $0.5$ steps over 100 consecutive episodes, typically near episode 750.

\paragraph{Greedy-GOMDP.} The risk-weighted greedy policy (Algorithm~\ref{alg:coord}) sorts sectors by belief-derived risk and assigns UAVs by proximity within energy budgets at $O(NZ)$ per step, providing a training-free baseline in the same environment.

\begin{algorithm}[!t]
\caption{Risk-Weighted Greedy Coordination}
\label{alg:coord}
\begin{algorithmic}[1]
\REQUIRE Belief $b_t$; fleet $\calU$; risk map $R_t$; sectors $\mathcal{P}$, $Z=|\mathcal{P}|=25$
\ENSURE Redeployment commands $\{(u_i,p_i)\}$
\STATE Compute $\{R^{p_i}_t\}$ from $b_t$; sort sectors descending by $R^{p_i}_t$
\FOR{each $u_i\in\calU$ (sorted by proximity)}
\STATE Assign $u_i$ to highest unassigned sector $p^*$ within energy budget
\ENDFOR
\STATE \textbf{return} commands \hfill $\triangleright$ $O(NZ)$ per step
\end{algorithmic}
\end{algorithm}

\paragraph{Two-stage verification and the operator model.} Stage-1 fusion computes $\Conf_t^{(1)}=w_H\hat{H}_t+w_W\hat{W}_t$ with $(w_H,w_W)=(0.65,0.35)$, and stage-2 applies a Bayesian update on secondary UAV verification, $\Conf_t^{(2)}=\Conf_t^{(1)}\cdot P(V_t\mid\mathrm{fire})/P(V_t)$ with $P(V_t\mid\mathrm{fire})=0.85$. Review is triggered if and only if $\Conf_t^{(2)}>\tau_2$, so every event an operator sees has cleared the threshold. We model the decision as a constant approval probability over that population, $\HA_t\sim\mathrm{Bernoulli}(1-p_{\mathrm{err}})$ with $p_{\mathrm{err}}=0.05$ and response delay $\mathcal{N}(3.0,0.8)$ steps truncated at zero. Both parameters are \emph{assumed}, not calibrated: no human-subject data supports them; the qualitative motivation is the oversight literature~\citep{holzinger2025}; and Appendix~\appSensitivity{} sweeps $p_{\mathrm{err}}$ to $0.20$.

\paragraph{Contract.} The chaincode logic sets $\mathrm{Alert}_t\leftarrow 1$ if and only if $\Conf_t^{(2)}>\tau$ and the Ed25519 signature verifies. Approved alerts generate a geofenced payload following the Common Alerting Protocol~\citep{oasis2010cap}; all events are committed with SHA-3-hashed evidence, and per-event nonces prevent replays. The named-authorizer role maps onto the incident-commander delegation chain of NIMS~\citep{fema2017nims}, a mapping we have not validated with a fire agency.

\section{Experimental Evaluation}
\label{sec:experiments}

\subsection{Setup}
\label{sec:setup}

\paragraph{Environment.} Wildfire dynamics unfold on a $100\times100$ grid with stochastic cellular-automaton propagation $P_{\mathrm{spread}}=\sigma(\alpha_1 W+\alpha_2 F-\alpha_3 H)$, weights $(0.45,0.35,0.30)$, tuned to 1--4 cells per 10-step window; $T=3000$ steps of 10\,s each. The long horizon is deliberate: detection occupies the first $0.5\%$ of an episode, and the remaining span generates the broadcast stream in the $F_p$ denominator of~\eqref{eq:fp}, so shortening it would make $F_p$ an artifact of episode length.

\paragraph{VIIRS-derived events.} Three VIIRS 375\,m active-fire events~\citep{schroeder2014} were preprocessed into $100\times100$ grids: California August Complex 2020, Evia 2021, and NSW Black Summer 2019--20. Hotspots were obtained from NASA FIRMS~\citep{firms2024} with ERA5 forcing~\citep{hersbach2020}; bounding boxes, dates, and hotspot counts are in the code package. Policies trained on the synthetic environment are evaluated zero-shot, and each event is mapped onto the same grid regardless of true extent, limiting what the comparison establishes about geographic scale.

\paragraph{Configurations.} Governed: \textbf{PPO-GOMDP}, \textbf{Greedy-GOMDP}. Enforcement without cryptography: \textbf{Shield-PPO} (same predicate, local logical shield after \citealp{alshiekh2018}) and \textbf{SafeLayer}~\citep{dalal2018}. Authentication without consensus: \textbf{Central+Sig} (BFT replaced by one Ed25519-verifying server). Policy-level surrogates: \textbf{PPO-CMDP}, \textbf{WCSAC}~\citep{yang2021wcsac}. Ungoverned: \textbf{Adaptive AI} (Greedy-GOMDP's coordination and verification, never gated, isolating the gate) and \textbf{Static}.

\paragraph{Seed protocol, statistics, and compute.} Each learned method is trained from 5 independent random initializations, and each trained policy is evaluated on 4 held-out environment seeds, yielding 20 paired evaluation runs per method over evaluation seeds 0--19; $\pm$ values are standard deviations across those runs, capturing both training and environment variance. Threshold selection and reward normalization use validation seeds 100--109. Significance tests are paired two-sided $t$-tests over the 20 runs with Holm--Bonferroni correction across a family of 8 hypotheses (PPO-GOMDP against each other configuration on $L_d$); all corrected $p<0.01$ unless noted. Compliance is definitional for the governed and shielded rows and undefined for the ungoverned rows, which attempt no authorization, so only the policy-level rows carry a measured value. Training one policy took $\approx 2.5$ GPU-hours on an NVIDIA RTX~3090 (24\,GB), AMD Ryzen~9 5950X, 64\,GB RAM, Ubuntu~22.04, PyTorch~2.1; the fleet sweep fine-tunes from $N{=}20$ checkpoints for 200 episodes. Cumulative budget: $\approx 110$ GPU-hours; inference: under 5\,ms on CPU.

\paragraph{Baseline scope.} PPO-CMDP, WCSAC, and SafeLayer are \emph{surrogates} reproducing the soft-constraint enforcement semantics of Lagrangian and projection methods, not reimplementations of the published algorithms, so their compliance figures characterize in-expectation enforcement in our environment rather than those methods at large. No published safe-RL implementation is run here, and a controlled comparison against public benchmark implementations, including SafeDreamer~\citep{huang2024safedreamer} and CCPO~\citep{yao2023ccpo}, is the highest-priority next step.

\subsection{Policy Comparison}
\label{sec:simresults}

Table~\ref{tab:rl_comparison} reports the main comparison. PPO-GOMDP reduces detection latency by 17.5\% relative to the training-free greedy baseline ($18.3\to15.1$ steps) at full authorization compliance. Against the \emph{trained} comparators, it is the slowest of the three: PPO-CMDP reaches $14.8$ steps and WCSAC $14.6$, but they violate the predicate in 7.2\% and 9.4\% of episodes, respectively. SafeLayer's learned projection reduces violations to 1.6\%, better than Lagrangian relaxation yet still short of a guarantee, and Shield-PPO matches GOMDP compliance in the benign setting while offering no defense against external injection.

\begin{table}[!t]
\caption{Policy comparison ($N{=}20$, 20 paired runs; mean $\pm$ std). $^{\mathrm{def}}$: compliance is definitional, not independently measured. N/A: the configuration attempts no authorization, so compliance is undefined. For governed rows, $\Le=L_d+4.2$ steps, a $+27.8\%$ overhead on $\Le$.}
\label{tab:rl_comparison}
\centering
\setlength{\tabcolsep}{2pt}
\scriptsize
\begin{tabular}{@{}lrrrrr@{}}
\toprule
\textbf{Method} & $L_d$ & $F_p$ (\%) & $FN_r$ (\%) & \textbf{Compl.} & \textbf{Enforce.} \\
\midrule
PPO-GOMDP & $15.1\!\pm\!1.1$ & $6.0\!\pm\!1.1$ & $2.1\!\pm\!0.9$ & $\mathbf{100.0\%}^{\mathrm{def}}$ & Crypto \\
Greedy-GOMDP & $18.3\!\pm\!1.4$ & $6.1\!\pm\!1.3$ & $2.3\!\pm\!1.0$ & $\mathbf{100.0\%}^{\mathrm{def}}$ & Crypto \\
Central+Sig & $15.0\!\pm\!1.1$ & $6.0\!\pm\!1.2$ & $2.1\!\pm\!0.9$ & $\mathbf{100.0\%}^{\mathrm{def}}$ & Sig.\ only \\
Shield-PPO & $15.2\!\pm\!1.2$ & $6.2\!\pm\!1.3$ & $2.2\!\pm\!0.9$ & $100.0\%^{\mathrm{def}}$ & Logical \\
SafeLayer & $14.9\!\pm\!1.1$ & $7.0\!\pm\!1.6$ & $2.4\!\pm\!1.0$ & $98.4\%$ & Learned \\
PPO-CMDP & $14.8\!\pm\!1.0$ & $8.3\!\pm\!2.4$ & $2.6\!\pm\!1.0$ & $92.8\%$ & Lagrangian \\
WCSAC & $14.6\!\pm\!1.2$ & $9.4\!\pm\!2.0$ & $3.8\!\pm\!1.4$ & $90.6\%$ & Lagrangian \\
Adaptive AI & $16.2\!\pm\!1.2$ & $22.4\!\pm\!2.1$ & $0.9\!\pm\!0.5$ & N/A & None \\
Static & $41.5\!\pm\!3.1$ & $15.3\!\pm\!2.4$ & $1.8\!\pm\!0.8$ & N/A & None \\
\bottomrule
\end{tabular}
\end{table}

A paired $t$-test on $L_d$ (PPO-GOMDP against PPO-CMDP) yields $t(19)=1.28$, two-sided $p=0.216$, which is non-significant but does not establish equivalence. We therefore apply two one-sided tests~\citep{lakens2017equivalence} with a pre-specified margin $\Delta=1.0$ step, or 10\,s of wall-clock time, which we treat as negligible relative to the pre-ignition window; this margin is a modeling judgment, not a field-calibrated threshold, since our fire model is unvalidated against operational spread data. With paired difference $0.3$ steps (SE $0.234$), $t_L(19)=5.55$ ($p<0.001$) and $t_U(19)=-2.99$ ($p=0.004$), so $\max(p_L,p_U)=0.004<0.05$ confirms equivalence within $\Delta$. Fleet-size trends across $N\in\{5,10,20,40\}$ preserve the ranking (Appendix~\appNotation).

\subsection{Ablation: Authentication versus Consensus}
\label{sec:ablation}

Table~\ref{tab:ablation} isolates each component and separates authentication from consensus. Injection resistance is measured rather than asserted: each attempt constructs a forged authorization and submits it through the contract entry point that a legitimate alert traverses, in three variants (unsigned certificate, valid signature from a non-validator key, and replay of an approved certificate). The outcome is deterministic for a correct implementation, so 100/100 confirms that the code path behaves as specified rather than estimating a rate.

\begin{table}[!t]
\caption{Ablation ($N{=}20$, 20 paired runs). Inj.: injection attacks blocked of 100 attempts, a deterministic check, not a sampled estimate. ``$-$Consensus'' keeps Ed25519 verification at one server; ``$-$All authentication'' removes signatures, nonces, and consensus.}
\label{tab:ablation}
\centering
\setlength{\tabcolsep}{2.5pt}
\footnotesize
\begin{tabular}{@{}lccr@{}}
\toprule
\textbf{Configuration} & $L_d$ & $F_p$ (\%) & \textbf{Inj.} \\
\midrule
PPO-GOMDP (full) & $15.1\!\pm\!1.1$ & $6.0\!\pm\!1.1$ & 100/100 \\
Greedy-GOMDP (full) & $18.3\!\pm\!1.4$ & $6.1\!\pm\!1.3$ & 100/100 \\
\;\,$-$ Adaptive coordination & $29.7\!\pm\!2.6$ & $6.1\!\pm\!1.2$ & 100/100 \\
\;\,$-$ HITL authorization & $15.2\!\pm\!1.1$ & $22.4\!\pm\!2.2$ & 100/100 \\
\;\,$-$ Consensus (Central+Sig) & $15.0\!\pm\!1.1$ & $6.0\!\pm\!1.2$ & 100/100 \\
\;\,$-$ All authentication & $15.1\!\pm\!1.1$ & $6.9\!\pm\!1.4$ & 0/100 \\
\;\,$-$ Multi-stage verification & $15.0\!\pm\!1.1$ & $14.8\!\pm\!2.0$ & 100/100 \\
PPO-CMDP (no contract) & $14.8\!\pm\!1.0$ & $8.3\!\pm\!2.4$ & 0/100 \\
\bottomrule
\end{tabular}
\end{table}

The decisive rows are the two authentication rows, which a single ``blockchain removed'' ablation would confound. Removing consensus while retaining signature verification changes nothing: Central+Sig still blocks all 100 injection attempts and matches PPO-GOMDP on $L_d$ and $F_p$ within measurement noise. Only removing authentication entirely admits injections. Signature verification alone is therefore sufficient against direct injection, and the consensus layer contributes no measurable robustness in any experiment we ran. The residual $F_p$ increase in the ``$-$All authentication'' row, from 6.0\% to 6.9\%, arises from losing per-event nonces: duplicate requests on already-authorized events are rebroadcast, and duplicates arriving after containment count as unwarranted. Ablating HITL raises $F_p$ to 22.4\%, matching ungoverned Adaptive AI, because both reduce to ungated two-stage verification. A lighter $m$-of-$n$ quorum (Appendix~\appMultisig) also matches on every measured metric.

\subsection{Adversarial Robustness}
\label{sec:robustness}

Table~\ref{tab:adversarial} reports results for all five declared threats. Spoofing raises $F_p$ everywhere because it corrupts the belief state directly and bypasses governance, but the HITL gate bounds the damage: at i.i.d.\ $p=0.20$, GOMDP reaches $9.4\%$ versus $38.7\%$ unauthenticated, and targeting high-belief cells at $p=0.10$ raises $F_p$ to $8.6\%$, above the $7.8\%$ under i.i.d.\ spoofing at the same rate. Tampering behaves similarly. Denial of service degrades latency and recall rather than precision, since dropped packets delay detection but cannot manufacture an authorization; at $p_{\mathrm{drop}}=0.40$, $L_d$ rises to $24.3$ steps and $FN_r$ to $6.1\%$. Compliance remains intact throughout, as no such attack can produce a valid certificate. Under validator compromise, the deterministic model confirms the tolerance boundary exactly: $f_c=2$ is tolerated and $f_c=3$ is not (Appendix~\appByzantine).

\begin{table}[!t]
\caption{Adversarial robustness ($N{=}20$, 20 paired runs): $F_p$ (\%) per attack, with $L_d$ for DoS. Central $=$ Adaptive AI without authentication.}
\label{tab:adversarial}
\centering
\setlength{\tabcolsep}{2pt}
\footnotesize
\begin{tabular}{@{}lcccc@{}}
\toprule
\textbf{Attack} & \textbf{Param.} & \textbf{GOMDP} & \textbf{Cent.+Sig} & \textbf{Central} \\
\midrule
No attack & --- & 6.0 & 6.0 & 22.4 \\
Spoofing (i.i.d.) & $p{=}0.05$ & 6.7 & 6.7 & 26.8 \\
Spoofing (i.i.d.) & $p{=}0.10$ & 7.8 & 7.9 & 31.2 \\
Spoofing (i.i.d.) & $p{=}0.20$ & 9.4 & 9.5 & 38.7 \\
Spoofing (strategic) & $p{=}0.10$ & 8.6 & 8.7 & 34.5 \\
Tampering & $p{=}0.10$ & 8.1 & 8.2 & 33.0 \\
Tampering & $p{=}0.20$ & 10.2 & 10.3 & 41.4 \\
DoS: $F_p$/$L_d$ & $p_{\mathrm{drop}}{=}0.20$ & 6.4/18.6 & 6.4/18.5 & 23.1/19.8 \\
DoS: $F_p$/$L_d$ & $p_{\mathrm{drop}}{=}0.40$ & 6.9/24.3 & 6.9/24.1 & 24.0/26.2 \\
Injection (blocked) & $p_{\mathrm{att}}{=}1$ & 100/100 & 100/100 & 0/100 \\
\bottomrule
\end{tabular}
\end{table}

\subsection{VIIRS-Derived Events}
\label{sec:realworld}

Zero-shot transfer to the three VIIRS-derived events preserves the ranking with degraded absolute performance (Appendix~\appVIIRS). Latency rises to 21.8--24.1 steps compared with 15.1 in the synthetic setting, reflecting sensor noise and geographic complexity, and $F_p$ rises to 7.9--9.1\% due to non-fire heat sources. The invariant holds in all episodes, and the surrogates' compliance deficit persists at 92.4--93.6\%. Adaptive AI achieves lower latency at 23.9--26.1\% false alert rates, confirming that unrestrained dissemination is hazardous under realistic noise. We do not report the IoT pipeline of \citet{tarigan2026} as a controlled improvement, since it differs in hardware, scenarios, modality, and latency definition.

\section{Discussion}
\label{sec:discussion}

\paragraph{What each layer buys.} Signature verification defeats direct injection and suffices for it; two-stage verification controls $F_p$ in the absence of attack, reducing it from 14.8\% to 6.0\%; the HITL gate separates 6.0\% from 22.4\%; and adaptive coordination separates $L_d=15.1$ from $29.7$. Byzantine consensus contributes nothing measurable to any of these. Its defensible justification is attributable non-repudiation across independently operated validators, with no single point of trust where no mutually trusted server exists. That is an architectural argument, not an empirical result, and Proposition~\ref{prop:correlated} shows that the analytic version does not survive correlated compromise against a hardened verifier. A trusted execution environment~\citep{costan2016} would plausibly match it at lower complexity; we did not evaluate one.

\paragraph{The full cost ledger.} The gate is evaluated after detection, so by construction it cannot change $L_d$: a policy detects a fire equally fast whether or not its alert will require authorization. What governance costs is the detection-to-broadcast delay, $\Delta_{\mathrm{BC}}+\Delta_{\mathrm{HITL}}\approx 4.2$ steps or 42\,s, a $27.8\%$ increase in $\Le$. The blockchain delay $\mathcal{N}(1.2,0.3)$ is assumed with no measurement behind it, and since the cost argument depends on it, Appendix~\appSensitivity{} reports the sensitivity. The second cost is more serious and cannot be omitted from a paper premised on missed detections being fatal. The gate raises $FN_r$ from 0.9\% to 2.1\%, a $2.3\times$ increase, and under plausible fatigue ($p_{\mathrm{err}}=0.20$) to 7.8\%, an $8.7\times$ increase. Every one of those missed detections represents a real fire whose alert the environment blocked. The mechanism introduces a new failure pathway: if the duty officer is unavailable during a mass-casualty event, a communications failure, or simultaneous incidents, the environment blocks a true alert \emph{by construction}---precisely in the high-load regime where a gate matters most. A deployment must weigh false alerts prevented against true alerts lost and specify a fallback authorization path; an $m$-of-$n$ quorum (Appendix~\appMultisig) improves availability, but we evaluated it only for injection resistance.

\paragraph{Limitations.} Six limitations bear on the safety case. First, Theorem~\ref{thm:safety} is conditional on Assumption~\ref{ass:impl}: the chaincode is not formally verified, so the guarantee pertains to the specification; applying existing verification machinery~\citep{bhargavan2016} is the highest-priority next step. Second, nothing is deployed: signing, key authorization, the nonce ledger, and threshold logic are real Python exercised by the reported attacks, but BFT ordering is a delay model, so consensus latency and throughput are unmeasured. Third, the fire model is a cellular automaton rather than a validated simulator such as FARSITE~\citep{finney1998}, so absolute latencies are within-model and $F_p$ is a within-model false-discovery ratio inflated by repeat alerts. Fourth, every comparator is ours, so the comparison bounds a regime rather than benchmarking published methods. Fifth, the operator model is a stationary Bernoulli approval with assumed parameters, omitting correlated fatigue and social engineering. Sixth, all three VIIRS events are forced onto a single grid, so nothing tests geographic scaling.

\section{Conclusion}
\label{sec:conclusion}

We introduced the Governance-Invariant MDP, in which the authorization constraint is a cryptographic state-transition invariant rather than a policy penalty, and stated its guarantee at the strength it holds: policy-agnostic under stated assumptions and dominated in deployment by validator compromise rather than cryptographic strength. The empirical contribution is a layered decomposition that prices each enforcement layer, including a negative result about the layer that motivated the architecture: signature-based non-repudiation makes environment-level enforcement robust to external injection, while Byzantine consensus adds nothing measurable and its analytic advantage disappears under correlated compromise against a hardened verifier. Formal chaincode verification, a validated fire simulator, a correlated-error operator model, and an availability-preserving quorum design are the prerequisites for the next stage.

\section*{Broader Impact and Path to Deployment}
\label{sec:broader_impact}

GOMDP addresses a civic need in which false alarms and missed detections both carry real social cost. Its contribution is accountability: every public-alert decision is authorized by a named operator and tamper-evidently logged, given faithful implementation and an honest validator supermajority. That accountability is not free, and the tradeoff belongs before a deploying agency, not in an appendix: the gate raises missed detections $2.3\times$ under our nominal operator model and $8.7\times$ under fatigue, so trading precision for recall may cost lives, and any fielded system needs a fallback when no authorizer is reachable. Field trials with fire-agency operators remain future work, UAV monitoring raises airspace and privacy concerns, and no component adds offensive capability.

\bibliography{references}

\end{document}
